\section{Experimental Setup}
\label{sec:setup}

\subsection{Study Design}
We use a full factorial design with the independent variables
\emph{LLM} (\TBD{} levels), \emph{programming language} (Python, Java, C++), and
\emph{hardware platform} (PC, server, Raspberry~Pi). Programming tasks are treated as a
random factor. The dependent variables are energy consumption (J, primary), execution
time, peak memory usage, and functional correctness.

\subsection{Large Language Models}
\TODO{Select models. Criteria: mix of proprietary and open-weight models, different
sizes/vendors, availability of a fixed model version/snapshot for reproducibility.
Record exact model identifiers, access dates, and decoding parameters.}

\begin{table}[h]
  \centering
  \caption{LLMs under study.}
  \label{tab:llms}
  \small
  \begin{tabular}{@{}llll@{}}
    \toprule
    Model & Vendor & Version / snapshot & Access \\
    \midrule
    \TBD & \TBD & \TBD & API / local \\
    \bottomrule
  \end{tabular}
\end{table}

\subsection{Programming Tasks}
\TODO{Select 10--15 tasks covering sorting, searching, numerical computation,
string processing, recursion, data structures, and memory-intensive workloads.
Requirements: (i) language-independent specification, (ii) a test suite for correctness,
(iii) a scalable input size so that each run lasts long enough for reliable energy
measurement (target: $\geq$\TBD{}\,s per run), (iv) low risk of training-data
contamination. Candidate sources: Computer Language Benchmarks Game, EvoEval,
EvalPerf, adapted LeetCode/Rosetta Code problems.}

\subsection{Hardware Platforms and Software Environment}

\begin{table}[h]
  \centering
  \caption{Hardware and software configuration.}
  \label{tab:platforms}
  \small
  \begin{tabular}{@{}llll@{}}
    \toprule
    & Desktop PC & Server & Raspberry~Pi \\
    \midrule
    CPU / architecture  & \TBD & \TBD & \TBD{} (ARM) \\
    RAM                 & \TBD & \TBD & \TBD \\
    OS / kernel         & \TBD & \TBD & \TBD \\
    Python (interpreter)& \TBD & \TBD & \TBD \\
    Java (JDK / JVM)    & \TBD & \TBD & \TBD \\
    C++ compiler, flags & \TBD & \TBD & \TBD \\
    Energy measurement  & \TBD & \TBD & \TBD \\
    \bottomrule
  \end{tabular}
\end{table}

\TODO{Use identical language-runtime versions on all platforms where possible
(e.g., via a pinned container image or toolchain). Fix CPU frequency governor,
disable turbo/boost where feasible, and document all deviations.}

\subsection{Energy Measurement}
\TODO{Decide measurement instruments per platform. Key issue: RAPL (PC, server) measures
CPU package/DRAM, whereas an external meter (Raspberry~Pi) measures whole-board power.
Options: (a) use external whole-system meters on all three platforms for comparability,
optionally complemented by RAPL; (b) compare rankings \emph{within} each platform rather
than absolute Joules \emph{across} platforms. Subtract idle baseline power.}
